\chapter{项目资源与源代码说明}

本附录整理了本研究涉及的原始数据集、实验脚本以及相关的开源项目地址，旨在提高研究的可复现性。

\section{开源项目地址}

本研究主要基于 PostgreSQL 数据库的演化数据进行分析，相关项目主页如下：

\begin{itemize}
    \item \textbf{PostgreSQL 官方仓库}：\url{https://github.com/postgres/postgres} 
    \item \textbf{本研究实验仓库}：\url{https://github.com/Colleygf/PostgreSQL-commit-analyis}
\end{itemize}

\section{项目结构说明}

本项目包含三个核心模块，分别对应论文中的静态分析、动态验证与形式化证明部分：

\begin{table}[htbp]
    \centering
    \caption{项目源代码目录说明}
    \begin{tabular}{llp{7cm}}
        \hline
        \textbf{目录名称} & \textbf{对应章节} & \textbf{核心功能说明} \\ \hline
        \texttt{static\_analyze/} & 第三章  & 包含基于 AST 的补丁特征提取脚本 (\texttt{diff\_extract.py})。  \\
        \texttt{dynamic\_analyze/} & 第四章  & 包含 Python 复现脚本、测试用例及生成的运行时日志 (\texttt{.log})。 \\
        \texttt{z3\_solver/} & 第五章  & 包含基于 Z3 求解器的状态抽象模型与验证脚本 (\texttt{verify\_fix.py})。  \\ \hline
    \end{tabular}
\end{table}
